\documentclass[../DoAn.tex]{subfiles}
\begin{document}

Chương~2 đã xác định A* và PIBT là hai thuật toán cần triển khai, Unity Engine là
nền tảng phát triển, và TCP là giao thức tích hợp máy chủ ngoài. Chương này trình bày
cơ sở lý thuyết và lý do lựa chọn từng công nghệ.

% ============================================================
\section{Bài toán MAPF}
\label{section:3.1}

MAPF được định nghĩa trên đồ thị $G = (V, E)$ với $k$ agent. Mỗi agent $i$ có vị trí
xuất phát $s_i \in V$ và đích $g_i \in V$. Tại mỗi bước thời gian, agent di chuyển
sang ô lân cận hoặc đứng yên. Ràng buộc: (i)~\textit{vertex conflict} -- hai agent
không được ở cùng ô cùng bước; (ii)~\textit{edge conflict} -- hai agent không được
đổi chỗ cho nhau trong cùng một bước. Mục tiêu: tìm tập kế hoạch $\pi_1, \ldots, \pi_k$
tối thiểu hóa tổng chi phí (số bước di chuyển) và không vi phạm ràng buộc \cite{stern2019mapf}.

Trong bối cảnh game thời gian thực, bài toán được đơn giản hóa: agent không cần đạt
đích cố định mà liên tục tiến gần mục tiêu di động (Eagle Base), và phải tái lập kế
hoạch khi bản đồ thay đổi (chướng ngại mới xuất hiện, agent bị tiêu diệt).

\section{Thuật toán A*}
\label{section:3.2}

A* \cite{hart1968astar} tìm đường ngắn nhất trên đồ thị có trọng số bằng cách mở
rộng các nút theo hàm ước lượng $f(n) = g(n) + h(n)$, trong đó $g(n)$ là chi phí
thực từ điểm xuất phát đến nút $n$, và $h(n)$ là heuristic ước lượng chi phí từ $n$
đến đích. Khi $h(n)$ nhất quán (consistent), A* đảm bảo tìm được đường ngắn nhất.

Trong đồ án, heuristic Manhattan $h(n) = |x_n - x_g| + |y_n - y_g|$ được dùng cho
lưới 4 láng giềng (lên, xuống, trái, phải), đảm bảo tính nhất quán và phù hợp với
cấu trúc bản đồ dạng ô vuông. A* được chọn làm baseline vì: (i) đã được kiểm chứng
rộng rãi trong game 2D; (ii) hiệu quả trên lưới thưa; (iii) dễ tích hợp và hiểu kết quả.

Hạn chế của A* trong bối cảnh MAPF: mỗi agent tính đường độc lập, không có cơ chế
phối hợp, dẫn đến va chạm và replan phản ứng liên tục.

\section{Thuật toán PIBT}
\label{section:3.3}

PIBT (Priority Inheritance with Backtracking) \cite{okumura2022pibt} là thuật toán MAPF
chạy theo bước thời gian (time-step based). Tại mỗi bước, các agent được xếp hạng ưu
tiên; agent ưu tiên cao hơn di chuyển trước, agent ưu tiên thấp hơn kế thừa ưu tiên
nếu bị chặn (\textit{priority inheritance}) và có thể hoàn tác bước di chuyển nếu
không tìm được ô trống (\textit{backtracking}).

Ưu điểm của PIBT: hoàn chỉnh trên bất kỳ bản đồ nào có đường đi tồn tại; thời gian
tính toán thực tế tỷ lệ tuyến tính với số agent trong trường hợp trung bình; phù hợp
với môi trường thời gian thực. PIBT được chọn vì nó cân bằng tốt giữa hiệu năng và
khả năng phối hợp đa agent, và đã được dùng thành công trong cuộc thi MAPF quốc tế
\cite{movingai2019benchmark}.

Đồ án triển khai PIBT theo hai cách: (i) trực tiếp trong Unity~C\# (\texttt{GridEnemyAgentPIBT})
để không phụ thuộc hạ tầng ngoài; (ii) qua giao tiếp TCP với máy chủ C++ tốc độ cao
(\texttt{GridEnemyAgentPIBT\_TCP}) để thể hiện hướng tích hợp planner bên ngoài Unity.

\section{Unity Engine và C\#}
\label{section:3.4}

Unity Engine \cite{unity2024engine} phiên bản 6000.3.10f1 (Unity~6) được lựa chọn vì:
(i) hỗ trợ xuất bản đa nền tảng bao gồm WebGL; (ii) hệ thống vật lý 2D (Rigidbody2D,
Collider2D) hoàn chỉnh; (iii) hệ sinh thái lớn và tài liệu đầy đủ; (iv) hỗ trợ
scripting~C\# với IL2CPP cho hiệu năng tốt. Unity NavMesh không được dùng làm hệ thống
điều hướng chính vì không tương thích với định dạng bản đồ \texttt{.map} và không
hỗ trợ ràng buộc MAPF ở lớp lập kế hoạch.

C\# được chọn (thay vì C++) vì tích hợp trực tiếp với Unity mà không cần lớp binding
bổ sung, giúp code dễ đọc và kiểm tra trong quá trình nghiên cứu.

Các lựa chọn thay thế đã xem xét: Godot Engine (thiếu hệ sinh thái C\# MAPF),
Unreal Engine (quá phức tạp cho 2D top-down), pygame (thiếu vật lý và mạng).

\section{Giao tiếp TCP giữa Unity và máy chủ C++}
\label{section:3.5}

Để tích hợp thuật toán PIBT viết bằng C++ (từ cuộc thi robot đua The League of Robot
Runners 2024) vào Unity mà không cần port toàn bộ sang C\#, đồ án sử dụng kết nối
TCP/IP theo mô hình client--server. Phía Unity đóng vai trò client: sau mỗi bước, Unity
gửi trạng thái hiện tại gồm vị trí các agent và kích thước bản đồ, rồi nhận lại ô mục
tiêu tiếp theo cho từng agent. Phía máy chủ C++ đóng vai trò server: chạy thuật toán
PIBT, duy trì trạng thái lập kế hoạch và lắng nghe trên cổng~7777 (localhost).

TCP được chọn thay vì UDP vì đảm bảo thứ tự và toàn vẹn gói tin -- yếu tố quan trọng khi
mỗi gói lệnh di chuyển đều cần thiết và không chấp nhận mất mát. Giao tiếp qua mạng nội
bộ (localhost/WSL) đảm bảo độ trễ thấp đủ cho ứng dụng thời gian thực.

\textbf{Kết chương~3:} Chương này đã trình bày cơ sở lý thuyết của MAPF, A* và PIBT,
đồng thời giải thích lý do lựa chọn từng công nghệ dựa trên yêu cầu xác định ở Chương~2.
Các nền tảng lý thuyết và kỹ thuật này là cơ sở để Chương~4 trình bày chi tiết thiết
kế và triển khai hệ thống.

\end{document}
